\subsection{Properties of Exponents} 
As well as addition and multiplication, we have a third ``operation'' that is (at the moment) designed only for some numbers.
\begin{definition}[Exponent]
For any real number $x$ and natural number $n$, the expression $x^n$ is defined as repeated multiplication.
\[
x^n=\underbrace{x\cdot x\cdots x}_{\text{$n$ copies}}
\]
The number $x$ is called the \term{base}, and the number $n$ is called the \term{exponent}.
\end{definition}
\begin{itemize}
\item We can read $x^2$ as ``$x$ \makeblank[1in]'' or ``$x$ to the second.''
\item We can read $x^3$ as ``$x$ \makeblank[1in]'' or ``$x$ to the third.''
\item For any other number $n$, we read $x^n$ as ``$x$ to the power $n$'' or ``$x$ to the $n$th power.''
\end{itemize}
We have certain properties of exponents, which are easy to see for natural number exponents.
\begin{displaybox}[Powers of the Same Base]
For any real number $x$ and natural numbers $m$, $n$, and $p$, with $m\geq n$, we have:
\[
x^m\cdot x^n = x^{m+n}
\qquad
\frac{x^m}{x^n}=x^{m-n}
\qquad
\frac{x^n}{x^m}=\frac{1}{x^{n-m}}
\qquad
\left(x^n\right)^p=x^{np}
\]
\end{displaybox}
\begin{Example}
Rewrite each expression with a single exponent.
\begin{lpic}{}{2}
\twocolleft{-1}{2}{$x^3\cdot x^5=$}{$\frac{t^5}{t^2}=$}
\twocollblleft{-2}{$\frac{3^{11}}{3^{29}}=$}{$\left(k^4\right)^5=$}
\end{lpic}
%\begin{itemize}
%\item $x^3\cdot x^5=\makeblank$
%\item $\dfrac{t^5}{t^2}=\makeblank$
%\item $\dfrac{3^{11}}{3^{29}}=\makeblank$
%\item $\left(k^4\right)^5=\makeblank$
%\end{itemize}
\end{Example}
\begin{displaybox}[Powers of Products and Quotients]
For any real numbers $x$, $y$, and $z$, with $z\neq 0$, and natural number $n$, we have:
\[
\left(xy\right)^n=x^ny^n
\qquad
\left(\frac{x}{z}\right)^n=\frac{x^n}{z^n}
\]
\end{displaybox}
\begin{Example}
Rewrite each expression with a single exponent.
\begin{lpic}{}{1}
\twocolleft{-1}{1}{$\left(2x\right)^3=$}{$\left(\dfrac{x}{6}\right)^2=$}
\end{lpic}
\end{Example}